\chapter{COMPLEX NUMBERS; MEASUREMENT OF ANGLES.}

We will not take up again here the complete exposé of the historical development of the theory of complex numbers or that of quaternions, these theories being essentially in the domain of Algebra (see pp.~72 and 62); but we will say a few words about the geometric representation of imaginary numbers which in many respects constitutes decisive progress in the history of Mathematics.

It is to C. F. Gauss that is due without dispute the first clear conception of the bijective correspondence between complex numbers and points of the plane, \(^{1}\) and especially the merit of having been the first to have known how to apply this idea to the theory of complex numbers, and to have foreseen clearly the advantage that the analysts of the XIXth century would make of it. During the course of the XVIIth and XVIIIth centuries, mathematicians had little by little reached the conviction that imaginary numbers, which allowed the solution of equations of the 2nd degree, also allowed the solution of algebraic equations of arbitrary degree. Numerous attempts at the proof of this theorem were published during the XVIIIth century; but, even without mentioning those that rely only on a vicious circle, there were none that were not open to serious objections. Gauss, after a detailed examination of these attempts and a tight critique of the gaps, proposed, in his inaugural Dissertation (written in 1797, appeared in 1799), to give finally a rigorous proof; taking up an idea given out in passing by d'Alembert (in the published proof by this latter in 1746, \(^{2}\) he remarks that the points \((a,b)\) of the plane such that \(a+ib\) is a root of the polynomial \(P(x+iy)=X(x,y)+iY(x,y)\) , are the intersections of the curves X=0 and Y=0; by a qualitative study of these curves, he shows that a continuous arc of one of them joins points of two distinct regions bounded by the other, and concludes from this that the curves meet ({[}124 a{]}, v. III, p.~3; see also {[}124 b{]}): a proof that, by its clarity and originality, constitutes a considerable progress on previous attempts, and is without doubt one of the first examples of an argument from pure Topology applied to a problem of Algebra. \(^{3}\)

In his Dissertation, Gauss does not define explicitly the correspondence between points of the plane and imaginary numbers; as far as these latter are concerned, and the questions about ``existence'' that they raised during two centuries, he even adopts a fairly reserved position, intentionally presenting all his arguments in a form in which only real numbers occur. But the march of ideas in his proof would be entirely unintelligible if they did not presuppose a fully conscious identification of the points of the plane and complex numbers; and his contemporaneous research on the theory of numbers and elliptic functions, where complex numbers also occur, only reinforces this hypothesis. At what point the geometric conception of imaginary numbers became familiar to him, and to what results it would lead in his hands, is what is shown clearly in the notes (published only in our day) where he applies complex numbers to the solution of problems of elementary Geometry ({[}124 a{]}, v. IV, p.~396 and v. VIII, p.~307). Still more explicit is the letter to Bessel of 1811 ({[}124 a{]}, v. VIII, pp.~90- 91), where he sketches the essentials of the theory of integration of functions of complex variables: ``In the same way'', he says, ``that one can represent the whole domain of real quantities by means of an indefinite straight line, so one can picture (``sinnlich machen'') the complete domain of all the quantities, the reals and the imaginaries, by means of an indefinite plane, where each point, determined by its abscissa a and its ordinate b, represents at the same time the quantity \(a + ib\) . The continuous movement of a value of x to another in consequence is made following a line, and so can be accomplished in an infinity of ways \ldots''

But it is only in 1831 that Gauss (in connection with the introduction of the ``Gaussian integers'' \(a + ib\) , where \(a\) and \(b\) are integers) displayed publicly his ideas on this point in such a clear a manner ({[}124 a{]}, v. II, Theoria Residuorum Biquadraticorum, Commentatio secunda, art. 38, p.~109, and Anzeige, pp.~174 ff.). In the interval, the idea of the geometric representation of imaginary numbers had been found again independently by two modest seekers, both amateur mathematicians, more or less self-taught, and for whom it was the only contribution to science, both also without much contact with the scientific environment of their time. Because of this, their work was in grave risk of being completely overlooked; it is precisely what happened for the first in time, the Dane C. Wessel, whose pamphlet, which appeared in 1798, very clearly conceived and drawn up, was only brought out of limbo a century later; and the same misadventure almost happened to the second, the Swiss J. Argand, who only by chance saw, in 1813, the exhumation of his work that he had published seven years earlier. \(^{4}\) This work provoked an active discussion in the Annales de Gergonne, and the matter was the object, in France and in England, of several publications (due to fairly obscure authors) between 1820 and 1830; but it lacked the authority of a big name to put an end to these controversies and to rally mathematicians to the new point of view; and it was necessary to wait until the middle of the century before the geometric representation of imaginary numbers was finally universally adopted, following the publications of Gauss (quoted above) in Germany, the work of Hamilton and Cayley on hypercomplex systems in England, and finally, in France the adherence of Cauchy, \(^{5}\) a few years only before Riemann, by a brilliant extension, came also to widen the role of Geometry in the theory of analytic functions, and create Topology at the same time.

The measurement of angles, by the arcs that they cut out on the circle, is as ancient as the notion of angle itself, and is already known to the Babylonians, from whom we have kept the unit of angle, the degree; it is besides only a matter for them of the measurement of angles between 0 and 360°, which sufficed for them, since their angles were mainly used by them to identify the position of celestial objects at predetermined points of their apparent trajectories, and to draw up tables for use in scientific and astrological situations.

With the Greek geometers of the classical era, the definition of angle (Eucl. El., I, def. 8 and 9) is still more restricted, since it only applies to angles less than two right angles; and as on the other hand their theory of ratios and of measurement relied on the comparison of arbitrarily large multiples of the quantities measured, the angles could not be for them a measurable quantity, although one finds naturally with them the notions of equal angles, of angles greater or less one than the other, and of the sum of two angles when this sum does not exceed two right angles. In the same way as for the addition of fractions, the measurement of angles must therefore have been in their eyes an empirical procedure without scientific value. This point of view is well illustrated by the admirable memoir of Archimedes on spirals ({[}5 b{]}, v. II, pp.~1-12), where, in default of being able to define these by the proportionality of the radius vector to the angle, he gives a cinematic description of them (def 1, p.~44; cf the statement of prop. 12, p.~46) from which he manages to draw out, as is shown by the rest of his work, all that the general notion of the measurement of angles would have given him if were to have had it. As for the Greek astronomers, they seem, on this point as on many others, to have been content to follow their Babylonian predecessors.

Here also, as in the evolution of the concept of real numbers (cf.~p.~150), the loosening of the spirit of rigour, during the decadence of Greek science, brought the return of the ``naive'' point of view, which in certain respects, is closer to our own than the rigid Euclidean conception. It is thus that an ill-advised interpolator inserts in Euclid the famous proposition (Eucl. El., VI, 33): ``The angles are proportional to the arcs that they cut out on the circle'', \(^{6}\) and an anonymous scholastic who comments on the ``proof'' of this proposition does not hesitate to introduce, without any justification of course, arcs equal to arbitrarily large multiples of a circumference, and the angles corresponding to these arcs ({[}107{]}, v. V, p.~357). But Vietus even, in the XVIth century, while appearing to touch on our modern conception of angle when he discovers that the equation \(\sin nx = \sin \alpha\) has several roots, only obtains the roots that correspond to angles less than 2 right angles ({[}319{]}, pp.~305-313). It is only in the XVIIth century that this point of view is overtaken in a definitive manner; and, after the discovery by Newton of the expansions in series of \(\sin x\) and \(\cos x\) had supplied expressions for these functions, valid for all values of the variables, one finds with Euler, in connection with logarithms of ``imaginary'' numbers, the precise conception of the notion of the measurement of an arbitrary angle ({[}108 a{]}, (1), v. XVII, p.~220).

Of course, the classical definition of the measurement of an angle by the length of the arc of a circle is not only intuitive, but essentially correct; however, it requires, in order to be made rigorous, the notion of the length of a curve, that is to say the integral Calculus. From the point of view of the structures that come into play, that is a very roundabout procedure, and it is possible to use no other methods than those of the theory of topological groups; the real exponential and the complex exponential appear thus as deriving from the same origin, the theorem characterising ``groups with one parameter''.

\chapter{METRIC SPACES.}

As we said (see p.~146), the notion of metric space was introduced in 1906 by M. Fréchet, and developed a few years later by F. Hausdorff in his ``Mengenlehre''. It acquired great importance after 1920, on the one hand following the fundamental work of S. Banach and his school on normed spaces and their applications to functional Analysis (see pp.~216 ff.), on the other by reason of the interest that was raised by the notion of valuation in Arithmetic and in Algebraic Geometry (where notably completion with respect to a valuation shows itself to be very fruitful).

From the period 1920-1930 are dated a whole series of studies undertaken by the Moscow school on the properties of the topology of a metric space, work which aimed in particular at obtaining necessary and sufficient conditions for a given topology to be metrizable. It is this movement of ideas that caused the appearance of interest in the notion of a normal space, defined in 1923 by Tietze, but whose important role was only recognised following the work of Urysohn {[}314{]} on the extension of continuous numerical functions. Apart from the trivial case of functions of a real variable, the problem of extension to the whole space of a continuous numerical function defined on a closed set, was treated for the first time (for the case of the plane) by H. Lebesgue {[}196 d{]}; before the definitive result of Urysohn, it had been solved for metric spaces by H. Tietze {[}307{]}. The extension of this problem to the case of functions with values in an arbitrary topological space acquired in these last years a considerable importance in algebraic Topology. Recent works have as well brought into the open that, in this type of question, the notion of normal space is hard to handle, because it still offers too many ``pathological'' possibilities; most of the time it is necessary to substitute for it the more restrictive notion of a paracompact space, introduced by J. Dieudonné in 1944 {[}90 a{]}; in this theory, the most remarkable result is the theorem due to A. H. Stone {[}300{]}, according to which every metrizable space is paracompact. \(^{1}\)

We have already focused attention (see p.~155) on the important works of the end of the XIXth century and of the beginning of the XXth century (E. Borel, Baire, Lebesgue, Osgood, W. H. Young) on the classification of sets of points in the spaces \(R^{n}\) , and on the classification and characterisation of numerical functions obtained by iterating, starting with continuous functions, the process of passing to the limit (for sequences of functions). It rapidly became apparent that metric spaces supplied a natural framework for research of this kind, whose development after 1910 is mainly due to the Russian and Polish schools. It is these schools that have between them brought to light the fundamental role played in modern Analysis by the notion of thin sets, and by the theorem on the countable intersection of open sets everywhere dense in a complete metric space, proved first (independently) by Osgood {[}240{]} for the numerical straight line and by Baire {[}11 b{]} for the spaces \(R^{n}\) .

On the other hand, in 1917, Souslin {[}289{]}, correcting an error of Lebesgue, showed that the continuous image of a Borel set is not necessarily Borel, which led him to the definition and study of the larger category of sets, since called ``analytic'' or ``Souslin''; after the premature death of Souslin this study was mainly undertaken by N. Lusin (whose ideas had inspired the work of Souslin) and by Polish mathematicians (see {[}209{]} and {[}189 b{]}). The actual importance of these sets relies mainly on their applications in the theory of integration (where, thanks to their special properties, they allow constructions that would be impossible on arbitrary measurable sets), and in modern potential theory, where the fundamental theorem on the capacity of Souslin sets, proved recently by G. Choquet {[}63{]}, has already shown itself to be rich in varied applications.
